\documentclass[11pt,a4paper]{article}
\usepackage[utf8]{inputenc}
\usepackage[T1]{fontenc}
\usepackage{amsmath,amssymb}
\usepackage[a4paper,margin=2.4cm]{geometry}
\usepackage{parskip}
\usepackage{titlesec}

\titleformat{\section}{\normalsize\bfseries}{\thesection.}{0.6em}{}
\allowdisplaybreaks

\newcommand{\Vs}{\mathcal{V}}
\newcommand{\Ws}{\mathcal{W}}
\newcommand{\Id}{\mathrm{Id}}
\newcommand{\cat}{\,\big\Vert\,}
\newcommand{\where}[1]{\par\nobreak\vspace{-2pt}{\small #1\par}\medskip}

\title{Multi-level NGP\,+\,GNN --- equations}
\author{}
\date{}

\begin{document}
\maketitle
\vspace{-2.2cm}

\section{The grid: levels, table length, finest cell}

\begin{equation}
L \;\in\; \mathbb{N} ,
\qquad
T \;\in\; \mathbb{N} ,
\qquad
p \;\geq\; 1
\label{eq:knobs}
\end{equation}
\where{where these are the three numbers the grid is described by:
$L$ is the number of levels;
$T$ sets the table length, $2^{T}$ rows, \emph{one table per level};
$p$ is the number of \emph{pixels covered by one cell of the finest level} --- $p = 1$ makes
the finest grid the pixel grid itself, and refining below the pixel spacing costs parameters
without moving the fit. Only two of $(L, p, b_d)$ are independent; which one is derived is
the choice made in \eqref{eq:bgiven}--\eqref{eq:bderived}.}

\begin{equation}
N_{\ell,d} \;=\; \max\Big(1,\; \min\big( \big\lfloor N_{\min,d}\, b_d^{\,\ell} \big\rceil ,\;
   N_{\max,d} \big)\Big) ,
\qquad \ell = 0,\dots,L
\label{eq:ladder}
\end{equation}
\where{where
$\ell \in \{0,\dots,L\}$ is the level index, $\ell = 0$ coarsest;
$d \in \{1,\dots,D\}$ is the axis index and $D$ the number of input dimensions
($D = 2$ for an image, $D = 4$ for $x,y,z,t$);
$N_{\ell,d} \in \mathbb{N}$ is the number of \emph{cells} along axis $d$ at level $\ell$,
hence $N_{\ell,d}+1$ vertices along that axis;
$N_{\min,d}$ is the coarsest resolution and $N_{\max,d}$ the cap, both per axis;
$b_d > 1$ is the per-axis growth factor of the \emph{resolution} per level, so the
\emph{cell width} is divided by $b_d$ at every step;
$\lfloor\cdot\rceil$ rounds to the nearest integer.}

\begin{align}
b_d &= \text{given}, \quad\text{default } 1.5 ,
&&N_{\max,d} = N_{\min,d}\, b_d^{\,L-1}
&&\text{(A) module}
\label{eq:bgiven}\\[4pt]
b_d &= \left( \frac{N_{\max,d}}{N_{\min,d}} \right)^{\!1/(L-1)} ,
&&N_{\max,d} = \max\Big( N_{\min,d}+1 ,\; \big\lfloor W_d / p \big\rceil \Big)
&&\text{(B) \texttt{gui\_image}}
\label{eq:bderived}
\end{align}
\where{where
$W_d$ is the extent of the data along axis $d$ \emph{in pixels} (or voxels, or microns ---
whatever unit $p$ is quoted in).}

\begin{align}
p_{\ell,d} &= W_d / N_{\ell,d}
&&\text{pixels per cell at level } \ell
\label{eq:pxlevel}\\[4pt]
\nu_\ell &= \prod_{d=1}^{D} \big(N_{\ell,d}+1\big)
&&\text{vertices at level } \ell
\label{eq:nverts}\\[4pt]
n_\ell &= \min\big(\nu_\ell,\; 2^{T}\big)
&&\text{rows level } \ell \text{ actually owns}
\label{eq:nrows}\\[4pt]
n_{\mathrm{enc}} &= F \sum_{\ell=0}^{L} \min\big(\nu_\ell,\; 2^{T}\big)
&&\text{encoder parameters}
\label{eq:nenc}
\end{align}
\where{where
$p_{\ell,d}$ is the number of pixels one level-$\ell$ cell covers along axis $d$, the number
worth watching level by level;
$\nu_\ell$ is how many vertices the level's lattice \emph{wants};
$2^{T}$ is how many rows it \emph{gets};
$n_{\mathrm{enc}}$ is the total encoder parameter count.}

\section{Prolongation and restriction}

\begin{align}
\boldsymbol\xi &= \mathbf{r}\odot \mathbf{N}_\ell - \big\lfloor \mathbf{r}\odot \mathbf{N}_\ell \big\rfloor ,
&&\mathbf{v} = \big\lfloor \mathbf{r}\odot \mathbf{N}_\ell \big\rfloor + \boldsymbol\sigma ,
\quad \boldsymbol\sigma \in \{0,1\}^{D}
\label{eq:cell}
\end{align}
\where{where
$\mathbf{r} \in [0,1]^{D}$ is the query point, normalised so the domain is the unit box;
$\mathbf{N}_\ell = (N_{\ell,1},\dots,N_{\ell,D})$ collects the level's resolutions;
$\odot$ is the elementwise product and $\lfloor\cdot\rfloor$ the floor, clamped so the last
cell of each axis is not overrun;
$\boldsymbol\xi \in [0,1)^{D}$ is the fractional position of $\mathbf{r}$ \emph{inside its
own level-$\ell$ cell};
$\boldsymbol\sigma \in \{0,1\}^{D}$ selects one of the $2^{D}$ corners of that cell;
$\mathbf{v} \in \mathbb{Z}^{D}$ is that corner's integer lattice coordinate.}

\begin{align}
w_\ell(\mathbf{v};\mathbf{r}) &= \prod_{d=1}^{D}
   \Big[\sigma_d\, s(\xi_d) + (1-\sigma_d)\big(1 - s(\xi_d)\big)\Big] ,
&&\sum_{\boldsymbol\sigma \in \{0,1\}^{D}} w_\ell = 1
\label{eq:weights}\\[4pt]
s(\xi) &= \xi
\quad\text{(linear)} ,
&& s(\xi) = \xi^2(3-2\xi)
\quad\text{(smoothstep)}
\label{eq:sfun}
\end{align}
\where{where
$w_\ell(\mathbf{v};\mathbf{r}) \in [0,1]$ is the interpolation weight of corner $\mathbf{v}$
for the query $\mathbf{r}$, and the weights of the $2^{D}$ corners sum to $1$ (partition of
unity), so interpolation reproduces constants exactly;
$s : [0,1] \to [0,1]$ is the per-axis interpolation profile --- linear makes the encoding
$C^{0}$, smoothstep makes it $C^{1}$ but destroys the nesting of \eqref{eq:V}, since
$s'(0) = s'(1) = 0$ forces zero slope at every vertex.}

\begin{align}
\big(I_\ell\, c\big)(\mathbf{r}) &= \sum_{\boldsymbol\sigma \in \{0,1\}^{D}}
   w_\ell(\mathbf{v};\mathbf{r})\; c_{\mathbf{v}}
&&\text{(prolongation)}
\label{eq:prolong}\\[4pt]
P_\ell &= I_\ell^{+} = \big(I_\ell^{\!\top} I_\ell\big)^{-1} I_\ell^{\!\top}
&&\text{(restriction)}
\label{eq:restrict}\\[4pt]
P_\ell I_\ell &= \Id_{n_\ell}
\label{eq:PI}
\end{align}
\where{where
$n_\ell = \min(\nu_\ell, 2^{T})$ of \eqref{eq:nrows} is the number of level-$\ell$ nodes --- grid
vertices here, or groups of units when the domain is a connectome;
$c \in \mathbb{R}^{n_\ell}$ holds one value per level-$\ell$ node;
$I_\ell : \mathbb{R}^{n_\ell} \to \mathbb{R}^{M}$ is the \emph{prolongation}, carrying node
values to the $M$ query points, and is exactly the encoding's own interpolation;
$I_\ell^{\!\top}$ is its adjoint, the backward scatter;
$I_\ell^{+}$ is the Moore--Penrose pseudo-inverse, i.e.\ the least-squares
\emph{restriction};
$\Id_{n_\ell}$ is the $n_\ell \times n_\ell$ identity.
Equation \eqref{eq:PI} is what the whole construction rests on: $P_\ell$ must be a true left
inverse. The cheaper mass-normalised adjoint $D_\ell^{-1}I_\ell^{\!\top}$ with
$D_\ell = \mathrm{diag}(I_\ell^{\!\top}\mathbf{1})$ does \emph{not} satisfy it.}

\begin{align}
R_\ell &= I_\ell P_\ell , & R_\ell^2 &= R_\ell , & R_{-1} &= 0
\label{eq:R}\\[2pt]
Q_\ell &= \Id - R_{\ell-1} = \Id - I_{\ell-1} P_{\ell-1} , & Q_0 &= \Id
\label{eq:Q}
\end{align}
\where{where
$R_\ell$ is the projector \emph{onto} the level-$\ell$ space --- it keeps whatever level
$\ell$ can represent and discards the rest --- and is idempotent by \eqref{eq:PI};
$Q_\ell$ is the \emph{band-pass}, which removes everything already representable at level
$\ell-1$, so that level $\ell$ can only contribute what is new;
$Q_0 = \Id$ because the coarsest level has nothing below it to subtract.}

\begin{align}
\Vs_\ell &= \operatorname{range}(I_\ell) ,
&\Vs_0 &\subset \Vs_1 \subset \cdots \subset \Vs_L
\label{eq:V}\\[2pt]
\Ws_\ell &= \Vs_\ell \ominus \Vs_{\ell-1} = Q_\ell \Vs_\ell = \Vs_\ell \cap \ker P_{\ell-1} ,
&\Vs_L &= \Ws_0 \oplus \Ws_1 \oplus \cdots \oplus \Ws_L
\label{eq:W}\\[2pt]
\dim \Ws_\ell &= n_\ell - n_{\ell-1} ,
&\sum_{\ell=0}^{L} \dim\Ws_\ell &= n_L
\label{eq:dimW}
\end{align}
\where{where
$\Vs_\ell \subset \mathbb{R}^{M}$ is everything level $\ell$ can express, of dimension
$n_\ell$;
the chain of inclusions is the \emph{nesting}, and it is exact only when each grid refines
the one above it --- integer $b_d$ and integer $N_{\min,d}$, per axis;
$\Ws_\ell$ is the \emph{detail space}, what level $\ell$ can express that level $\ell-1$
cannot;
$\ominus$ is the complement of $\Vs_{\ell-1}$ inside $\Vs_\ell$ taken along $\ker P_{\ell-1}$
(oblique, not orthogonal, unless $P_{\ell-1}$ is the orthogonal projection);
$\oplus$ is a direct sum, and it is the direct sum that makes each level's contribution
recoverable from the total;
$\dim\Ws_\ell$ telescoping to $n_L$ is the numerical test that the nesting is real.}

\section{The encoder}

\begin{equation}
\big(u_i,\; a_i\big) \;=\; \mathrm{NGP}\big(x_i,\, y_i,\, z_i,\, t_i\big) ,
\qquad i = 1,\dots,N
\label{eq:ngp}
\end{equation}
\begin{equation}
\mathbf{r}_i = (x_i,\,y_i,\,z_i,\,t_i) \in [0,1]^{4}
\label{eq:coord}
\end{equation}
\where{where
$i$ indexes a sample --- one neuron at one time, or one voxel at one time --- and $N$ is
their number;
$(x_i,y_i,z_i) \in \mathbb{R}^{3}$ is a position and $t_i \in \mathbb{R}$ a time, both
rescaled into the unit box as $\mathbf{r}_i$;
$u_i \in \mathbb{R}^{d_u}$ is the \emph{state} at that sample --- a voltage, a fluorescence
$\Delta F/F$ --- with $d_u = 1$ for a scalar recording;
$a_i \in \mathbb{R}^{d_a}$ is the \emph{embedding} at that sample, the object that replaces
the free per-neuron parameter, with $d_a$ its width;
$\mathrm{NGP}$ is the multiresolution hash encoding, expanded next.}

\section{The encoder expanded: hash table and MLP}

\begin{align}
h_\ell(\mathbf{v}) &=
\begin{cases}
\displaystyle \sum_{d=1}^{D} v_d \prod_{e<d} \big(N_{\ell,e}+1\big) ,
  & \displaystyle \prod_{d}\big(N_{\ell,d}+1\big) \le 2^{T}
  \quad\text{(dense)} \\[10pt]
\displaystyle \Big( \bigoplus_{d=1}^{D} v_d\, \pi_d \Big) \bmod 2^{T} ,
  & \text{otherwise} \quad\text{(hashed)}
\end{cases}
\label{eq:hash}\\[6pt]
\boldsymbol\pi &= \big(1,\; 2654435761,\; 805459861,\; 3674653429\big)
\label{eq:primes}
\end{align}
\where{where
$h_\ell(\mathbf{v}) \in \{0,\dots,2^{T}-1\}$ is the table row that vertex $\mathbf{v}$ reads at
level $\ell$;
$2^{T}$ is the number of rows available to one level, one table per level;
the top branch applies when the level's lattice fits in the table, so the map is a bijection
and there are \emph{no} collisions;
the bottom branch is Instant-NGP's spatial hash, where $\bigoplus$ is bitwise XOR and
$\pi_d$ are large primes chosen to scatter collisions pseudo-randomly across the domain, so
that two colliding vertices are typically far apart;
each level owns a \emph{disjoint} slice of the flat table, so competition for rows is
strictly within a level.}

\begin{align}
\mathbf{z}_i^{(\ell)} &= \sum_{\boldsymbol\sigma\in\{0,1\}^{D}}
   w_\ell(\mathbf{v};\mathbf{r}_i)\;
   \boldsymbol\theta^{(\ell)}_{\,h_\ell(\mathbf{v})}
\;\in\; \mathbb{R}^{F} ,
&&\boldsymbol\theta^{(\ell)} \in \mathbb{R}^{2^{T}\times F}
\label{eq:lookup}\\[4pt]
\mathbf{z}_i &= \Big[\, \mathbf{z}_i^{(0)} \cat \mathbf{z}_i^{(1)} \cat \cdots \cat \mathbf{z}_i^{(L)} \,\Big]
\;\in\; \mathbb{R}^{(L+1)F}
\label{eq:concat}\\[4pt]
\big(u_i,\; a_i\big) &= \mathrm{MLP}_{\psi}\big(\mathbf{z}_i\big)
\label{eq:mlp}\\[4pt]
a_i^{(\ell)} &\equiv \mathbf{z}_i^{(\ell)}
\label{eq:aslice}
\end{align}
\where{where
$\boldsymbol\theta^{(\ell)} \in \mathbb{R}^{2^{T} \times F}$ is the level-$\ell$ feature table,
the only place learned capacity is stored, initialised uniform on
$[-10^{-4}, 10^{-4}]$;
$F \in \mathbb{N}$ is the number of features per row (typically $2$ or $4$);
$\boldsymbol\theta^{(\ell)}_{k}$ is row $k$ of that table;
$\mathbf{z}_i^{(\ell)} \in \mathbb{R}^{F}$ is the interpolated level-$\ell$ feature at sample
$i$, i.e.\ $\eqref{eq:prolong}$ applied to the table;
$\cat$ denotes concatenation, so $\mathbf{z}_i$ stacks all $L+1$ levels;
$\mathrm{MLP}_{\psi}$ is the decoder with parameters $\psi$, mapping the concatenated
features to the state and the embedding;
\eqref{eq:aslice} is the alternative in which the embedding is \emph{not} decoded but is the
raw per-level slice, which is what makes $a$ multiresolution rather than a single vector.}

\section{Encoder and operator together, multi-level}

\begin{equation}
\partial_t u \;=\; \mathcal{F}\big[u\big]
\label{eq:pde}
\end{equation}
\where{where
$\mathcal{F}$ is the true (unknown) evolution operator --- the thing being recovered ---
mapping a state field to its rate of change.}

\begin{equation}
\boxed{\;
\dot u^{*}(\mathbf{r},t) \;=\; \sum_{\ell=0}^{L}
Q_\ell\, I_\ell \Big[\, \mathrm{GNN}_\ell\big(\, P_\ell u,\; a_\ell \,\big) \Big](\mathbf{r},t)
\;}
\label{eq:model}
\end{equation}
\where{where
$\dot u^{*} \in \mathbb{R}^{M \times d_u}$ is the \emph{predicted} rate of change at the
query points, the model's output;
$\mathrm{GNN}_\ell$ is the level-$\ell$ operator: it reads the restricted state and the
level's embedding at every level-$\ell$ node and emits one rate of change per node, so its
output is band-limited from \emph{above} by the level's own cell size;
$I_\ell$ carries that back to full resolution and $Q_\ell$ removes the part a coarser level
could already have produced, band-limiting from \emph{below};
the sum over $\ell$ is a sum of \emph{contributions}, not of independent predictions.}

\begin{align}
W^{(\ell)} &= P_\ell\, W\, I_\ell
\label{eq:coarseW}\\[4pt]
\mathbf{m}^{(\ell)}_\alpha &= \sum_{\beta} W^{(\ell)}_{\alpha\beta}\;
   \psi_\ell\big(\, v_\alpha,\, v_\beta,\, a_{\ell,\alpha},\, a_{\ell,\beta} \,\big)
\label{eq:message}\\[4pt]
\big[\mathrm{GNN}_\ell(v, a_\ell)\big]_\alpha &=
   \phi_\ell\big(\, v_\alpha,\; a_{\ell,\alpha},\; \mathbf{m}^{(\ell)}_\alpha \,\big) ,
\qquad v = P_\ell u
\label{eq:update}
\end{align}
\where{where
$W \in \mathbb{R}^{N \times N}$ is the connectivity at the finest scale --- the synaptic
weight from unit $\beta$ to unit $\alpha$, or lattice adjacency in the volumetric case;
$W^{(\ell)} \in \mathbb{R}^{n_\ell \times n_\ell}$ is its Galerkin coarsening, which for a
partition by cell type is exactly the type-to-type connectivity matrix;
$\alpha, \beta \in \{1,\dots,n_\ell\}$ index level-$\ell$ nodes;
$v = P_\ell u \in \mathbb{R}^{n_\ell \times d_u}$ is the restricted state, and $v_\alpha$ its
value at node $\alpha$;
$a_{\ell,\alpha} \in \mathbb{R}^{F}$ is that node's level-$\ell$ embedding;
$\psi_\ell$ is the message function, a small network evaluated once per edge;
$\mathbf{m}^{(\ell)}_\alpha$ is the aggregated message at node $\alpha$, a weighted sum over
its neighbours;
$\phi_\ell$ is the update function, producing that node's contribution to the rate of change.}

\begin{equation}
\mathcal{L} \;=\;
\sum_i \big\| u_i' - u_i \big\|^2
\;+\; \lambda \sum_i \big\| \dot u_i^{*} - \dot u_i \big\|^2
\label{eq:loss}
\end{equation}
\where{where
$u_i'$ is the encoder's reconstruction of the state and $u_i$ the observation, so the first
term is the \emph{data fit};
$\dot u_i$ is the observed rate of change --- a finite difference of the recording, or a
leave-one-out local-linear slope --- and the second term is the \emph{model fit};
$\lambda > 0$ balances them, and carries units of time$^2$ because the two terms differ by
$1/\text{time}^2$, so it does not transfer between datasets of different frame rate unless
both terms are non-dimensionalised first.}

\section{Everything expanded}

\begin{align}
\mathbf{z}_i^{(\ell)} &= \sum_{\boldsymbol\sigma\in\{0,1\}^{D}}
   w_\ell(\mathbf{v};\mathbf{r}_i)\; \boldsymbol\theta^{(\ell)}_{\,h_\ell(\mathbf{v})}
&&\text{table}
\label{eq:f-lookup}\\[4pt]
u_i' &= \mathrm{MLP}_{\psi}\Big(\big[\, \mathbf{z}_i^{(0)} \cat \cdots \cat \mathbf{z}_i^{(L)} \,\big]\Big)
&&\text{data head}
\label{eq:f-decode}\\[4pt]
a_{\ell,\alpha} &= \mathbf{z}_\alpha^{(\ell)} ,
\qquad v^{(\ell)} = P_\ell\, u
&&\text{level inputs}
\label{eq:f-inputs}\\[4pt]
\mathbf{m}^{(\ell)}_\alpha &= \sum_{\beta} \big(P_\ell W I_\ell\big)_{\alpha\beta}\,
   \psi_\ell\big(v^{(\ell)}_\alpha, v^{(\ell)}_\beta, a_{\ell,\alpha}, a_{\ell,\beta}\big)
&&\text{messages}
\label{eq:f-msg}\\[4pt]
g^{(\ell)}_\alpha &= \phi_\ell\big(v^{(\ell)}_\alpha,\; a_{\ell,\alpha},\; \mathbf{m}^{(\ell)}_\alpha\big)
&&\text{level output}
\label{eq:f-out}\\[4pt]
\dot u^{*} &= \sum_{\ell=0}^{L}
   \big(\Id - I_{\ell-1}P_{\ell-1}\big)\, I_\ell\, g^{(\ell)}
&&\text{band-pass synthesis}
\label{eq:f-synth}\\[4pt]
\mathcal{L} &= \sum_i \big\| u_i' - u_i \big\|^2
   \;+\; \lambda \sum_i \big\| \dot u_i^{*} - \dot u_i \big\|^2
&&\text{objective}
\label{eq:f-loss}
\end{align}
\where{where all symbols are as defined above, and
$g^{(\ell)} \in \mathbb{R}^{n_\ell \times d_u}$ is the level-$\ell$ output vector, one rate
of change per level-$\ell$ node, before it is carried back to full resolution;
$v^{(\ell)} = P_\ell u$ is the state restricted to level $\ell$, so the operator reads the
\emph{state} and the encoder supplies only the \emph{embedding}.}

\begin{equation}
\Theta \;=\; \Big\{\, \boldsymbol\theta^{(0)},\dots,\boldsymbol\theta^{(L)} ,\;
\psi ,\; \psi_0,\dots,\psi_L ,\; \phi_0,\dots,\phi_L \,\Big\}
\label{eq:params}
\end{equation}
\where{where
$\Theta$ is everything learned: the $L+1$ feature tables, the decoder parameters $\psi$, and
one message and one update network per level. Sharing $\psi_\ell$ and $\phi_\ell$ across
levels, conditioned on the level's cell size, replaces the last $2(L+1)$ networks by two.}

\begin{equation}
b_d \in \mathbb{N} ,
\qquad N_{\min,d} \in \mathbb{N} ,
\qquad s(\xi) = \xi ,
\qquad P_\ell I_\ell = \Id
\label{eq:conditions}
\end{equation}
\where{where these are the four conditions under which \eqref{eq:f-synth} is a genuine
band-pass rather than a relabelling: integer growth factor and integer coarsest resolution
\emph{per axis}, so the grids are exactly nested; linear interpolation on every band-passed
axis, since smoothstep breaks the nesting; and a restriction that is a true left inverse,
which requires solving the level-$\ell$ Gram system rather than scattering.}

\end{document}
